\begin{lemma}
\label{editorial-source-lemma-invert-closed-quotient}
Let $R$ be a ring. Let $S \subset R$ be a multiplicative subset.
Assume the image of the map $\Spec(S^{-1}R) \to \Spec(R)$
is closed. Then $S^{-1}R \cong R/I$ for some ideal $I \subset R$.
\end{lemma}

\begin{proof}
Let $I = \Ker(R \to S^{-1}R)$ so that $V(I)$ contains the image.
Say the image is the closed subset $V(I') \subset \Spec(R)$ for
some ideal $I' \subset R$. So $V(I') \subset V(I)$.
For $f \in I'$ we see that $f/1 \in S^{-1}R$
belongs to every prime ideal. Hence $f^n$ maps to zero in $S^{-1}R$
for some $n \geq 1$ (Lemma \ref{editorial-source-lemma-Zariski-topology}).
Hence $V(I') = V(I)$.
Every $g\in S$ avoids all primes in the image, hence every prime
containing $I$. Its class $\overline g$ in $R/I$ belongs to no maximal
ideal and is therefore a unit. The quotient and localization properties
give ring maps
$$
\alpha:R/I\longrightarrow S^{-1}R,\quad \overline r\longmapsto r/1,
\qquad
\beta:S^{-1}R\longrightarrow R/I,\quad r/s\longmapsto
\overline r\,\overline s^{-1}.
$$
For every $\overline r$, $\beta\alpha(\overline r)=\overline r$.
For every $r/s$, $\alpha\beta(r/s)=(r/1)(s/1)^{-1}=r/s$.
Thus the maps are inverse, and their composite with $R\to R/I$
is the original localization map. In particular the localization map
is surjective. Conversely, if $R\to S^{-1}R$ is surjective, it is
the quotient by its kernel and its spectrum has closed image by
Lemma \ref{algebra-lemma-spec-closed}.
\end{proof}